% ============================================================
%  PROYECTO DE INVESTIGACIÓN  —  MIA-303
%  Carlos Pérez Pérez  —  UNI  2026
%
%  Versión standalone: no requiere TesisUNI.cls
%  Usa clases y paquetes estándar de LaTeX.
% ============================================================

\documentclass[11pt,a4paper,oneside]{book}

% ── Codificación ──────────────────────────────────────────
\usepackage[utf8]{inputenc}
\usepackage[T1]{fontenc}

% ── Idioma ────────────────────────────────────────────────
\usepackage[main=spanish,english]{babel}

% ── Fuente Times New Roman ────────────────────────────────
\usepackage{mathptmx}       % Times en texto y fórmulas
\usepackage{textcomp}

% ── Márgenes (3 izq / 2.5 der / 2.5 sup / 2.5 inf) ──────
\usepackage[left=3cm, right=2.5cm, top=2.5cm, bottom=2.5cm]{geometry}

% ── Interlineado ──────────────────────────────────────────
\usepackage{setspace}

% ── Encabezados y pies ────────────────────────────────────
\usepackage{fancyhdr}
\pagestyle{fancy}
\fancyhf{}
\fancyhead[R]{\small\slshape\leftmark}
\fancyfoot[C]{\thepage}
\renewcommand{\headrulewidth}{0.4pt}

% ── Formato de capítulos y secciones ─────────────────────
\usepackage{titlesec}
% Capítulos: centrados, mayúsculas, sin "Capítulo N"
\titleformat{\chapter}[block]
  {\normalfont\large\bfseries\centering}
  {\thechapter.}
  {1em}
  {\MakeUppercase}
\titlespacing*{\chapter}{0pt}{20pt}{14pt}

% Secciones numeradas, negrita
\titleformat{\section}
  {\normalfont\normalsize\bfseries}
  {\thesection}
  {1em}
  {}
\titlespacing*{\section}{0pt}{12pt}{6pt}

% Subsecciones
\titleformat{\subsection}
  {\normalfont\normalsize\bfseries}
  {\thesubsection}
  {1em}
  {}
\titlespacing*{\subsection}{0pt}{10pt}{4pt}

% ── Matemáticas ───────────────────────────────────────────
\usepackage{amssymb,amsmath,amsbsy}
\usepackage{mathdots}
\usepackage{mathrsfs}

% ── Paquetes del template original ────────────────────────
\usepackage[table,xcdraw]{xcolor}
\usepackage{multirow}
\usepackage{tikz}
\usepackage[printonlyused]{acronym}
\usepackage{gensymb}
\usepackage{enumitem}

% ── Hipervínculos (PDF interactivo) ───────────────────────
\usepackage[hidelinks,
            pdftitle={Detección Automatizada de Eventos Adversos - GEMSES - MIMIC-IV},
            pdfauthor={Carlos Pérez Pérez}]{hyperref}

% ── Sangría y párrafos ────────────────────────────────────
\setlength{\parindent}{1.25cm}
\setlength{\parskip}{4pt}

% ── Compatibilidad con comandos de TesisUNI ───────────────
\newcommand{\BOthers}[1]{et al.\hbox{}}

\hyphenation{Pro-gram-ma-ble}

% ── Datos del proyecto ────────────────────────────────────
% ============================================================
%  DATOS GENERALES DEL PROYECTO
% ============================================================
% Este archivo puede ser utilizado para definir macros
% reutilizables en toda la tesis.

\newcommand{\tituloTesis}{%
    Detección Automatizada y Priorización de Eventos Adversos en
    Notas Clínicas mediante un Flujo de Procesamiento (Pipeline) de
    NLP y Aprendizaje Automático --- Aplicación de la Matriz de
    Priorización del Modelo GEMSES sobre MIMIC-IV%
}

\newcommand{\autorTesis}{Carlos Pérez Pérez}
\newcommand{\cursoTesis}{MIA-303 --- Proyectos de Investigación I}
\newcommand{\docenteTesis}{Mg. María Tejada Begazo}
\newcommand{\universidadTesis}{Universidad Nacional de Ingeniería}
\newcommand{\facultadTesis}{Facultad de Ingeniería Industrial y de Sistemas}
\newcommand{\programaTesis}{Maestría en Inteligencia Artificial}
\newcommand{\anioTesis}{2026}
\newcommand{\ciudadTesis}{Lima --- Perú}


% ============================================================
\begin{document}

\renewcommand{\BOthers}[1]{et al.\hbox{}}
\onehalfspacing

% ── PARTE INICIAL ──────────────────────────────────────────
\frontmatter

% ============================================================
%  CARÁTULA
% ============================================================

\begin{titlepage}
\centering

{\large\bfseries UNIVERSIDAD NACIONAL DE INGENIERÍA}\\[4pt]
{\normalsize\bfseries FACULTAD DE INGENIERÍA INDUSTRIAL Y DE SISTEMAS}\\[4pt]
{\normalsize\bfseries UNIDAD DE POSGRADO}\\[2pt]
{\normalsize\bfseries MAESTRÍA EN INTELIGENCIA ARTIFICIAL}\\[40pt]

{\Large\bfseries\color{black} PROYECTO DE INVESTIGACIÓN}\\[28pt]

\begin{minipage}{0.88\textwidth}
\centering
{\normalsize\bfseries
«DETECCIÓN AUTOMATIZADA Y PRIORIZACIÓN DE EVENTOS ADVERSOS EN
NOTAS CLÍNICAS MEDIANTE UN FLUJO DE PROCESAMIENTO (PIPELINE) DE
NLP Y APRENDIZAJE AUTOMÁTICO --- APLICACIÓN DE LA MATRIZ DE
PRIORIZACIÓN DEL MODELO GEMSES SOBRE MIMIC-IV»
}
\end{minipage}

\vspace{28pt}
{\normalsize\bfseries TRABAJO PARCIAL --- CAPÍTULOS I Y II}

\vfill

\begin{flushleft}
{\normalsize Elaborado por:}\\[4pt]
{\normalsize\bfseries Carlos Pérez Pérez}\\[14pt]
{\normalsize Curso: MIA-303 --- Proyectos de Investigación I}\\[6pt]
{\normalsize Docente: Mg. María Tejada Begazo}
\end{flushleft}

\vspace{30pt}
{\normalsize Lima --- Perú, 2026}

\end{titlepage}


% Índice
\singlespacing
\renewcommand\contentsname{\centering ÍNDICE}
\tableofcontents
\onehalfspacing

% ── PARTE CENTRAL ──────────────────────────────────────────
\mainmatter

% Capítulo I: Planteamiento del Problema y Justificación
% ============================================================
%  CAPÍTULO I: PLANTEAMIENTO DEL PROBLEMA Y JUSTIFICACIÓN
% ============================================================

\chapter{PLANTEAMIENTO DEL PROBLEMA Y JUSTIFICACIÓN}
\label{cap:problema}

% ─────────────────────────────────────────────────────────────
\section{Planteamiento del Problema}
\label{sec:planteamiento}

La seguridad del paciente constituye uno de los principales retos de los sistemas de salud contemporáneos. Según la Organización Mundial de la Salud \mbox{(OMS, 2019)}, los eventos adversos derivados de la atención sanitaria figuran entre las diez principales causas de muerte y discapacidad en el mundo, generando además costos sustanciales para los sistemas asistenciales. En el Perú, el Registro Nacional de Instituciones Prestadoras de Servicios de Salud (RENIPRESS) registra 26{,}663 establecimientos activos que atienden a más de 33 millones de personas, distribuidos en doce tipos institucionales: sector privado (16{,}685 IPRESS), Gobierno Regional (8{,}380), MINSA (510), EsSalud (408), Sanidades de las Fuerzas Armadas y Policiales, SISOL, gobiernos locales y otros. EsSalud, que concentra la mayor densidad de hospitalizaciones de tercer nivel, aprobó en 2021 la Directiva GG-ESSALUD «Registro, Notificación y Gestión de los Eventos Relacionados con la Seguridad del Paciente (ERSP)», la cual incorpora explícitamente la Matriz de Priorización del Modelo de Gestión Moderna de los Servicios de Salud (GEMSES) como herramienta normativa para clasificar el nivel de riesgo y de gestión de los eventos.

Sin embargo, el proceso de identificación y priorización depende todavía del registro voluntario de los profesionales de la salud y del análisis manual de las descripciones libres consignadas en las notas clínicas. Esta dependencia genera tres obstáculos persistentes: (i)~la subnotificación voluntaria, asociada al temor de repercusiones laborales; (ii)~la dependencia de descripciones en lenguaje natural que dificultan su codificación, agregación y análisis sistemático; y (iii)~la falta de mecanismos automatizados que prioricen los eventos de mayor riesgo para escalarlos al nivel de gestión adecuado en tiempo oportuno.

En el caso específico de EsSalud, la Directiva GG-ESSALUD-2021 establece un flujo formal de Registro $\rightarrow$ Notificación $\rightarrow$ Validación $\rightarrow$ Codificación $\rightarrow$ Priorización $\rightarrow$ Mitigación $\rightarrow$ Análisis de Causa Raíz $\rightarrow$ Plan de Acción. Aunque la directiva normaliza la Matriz de Priorización GEMSES como instrumento oficial, su aplicación efectiva requiere que personal de la Oficina de Gestión de la Calidad y Humanización procese manualmente cada notificación, lo que constituye un cuello de botella estructural que hace el proceso lento, costoso y vulnerable a la heterogeneidad en la interpretación.

La presente investigación propone aprovechar los avances recientes en Procesamiento de Lenguaje Natural (NLP) y aprendizaje automático para automatizar la detección de 188 de los 231 eventos adversos del Anexo 02 de la Directiva GG-ESSALUD-2021 en notas clínicas no estructuradas, y la posterior aplicación de la Matriz de Priorización GEMSES. Dada la ausencia de acceso a registros institucionales peruanos, el estudio se desarrollará sobre el \textit{dataset} abierto MIMIC-IV (\textit{Medical Information Mart for Intensive Care}, versión~IV). El estudio se concibe como un \textbf{demostrador técnico de viabilidad}: demuestra que la automatización GEMSES es realizable con datos clínicos reales, generando un protocolo documentado y replicable que las instituciones peruanas ---cuyo sistema nacional aún carece de soporte normativo e infraestructura para esta automatización--- podrán adoptar cuando se habiliten los accesos correspondientes.

% ─────────────────────────────────────────────────────────────
\section{Formulación del Problema}
\label{sec:formulacion}

La gestión de eventos adversos en los servicios de salud presenta brechas críticas en la detección oportuna, la codificación consistente y la priorización basada en evidencia. Los modelos manuales dependen del criterio subjetivo del profesional y generan subnotificación sistemática. Más aún, incluso cuando un evento se reporta, la cadena de gestión presenta cuatro puntos de quiebre sucesivos: (i)~el evento se registra pero \textit{no se analiza}, por saturación del personal de calidad; (ii)~se analiza pero \textit{no se prioriza}, por ausencia de herramientas de clasificación estandarizada; (iii)~se prioriza pero \textit{no se asigna un responsable} con plazo definido; y (iv)~se gestiona puntualmente pero \textit{no existe un indicador permanente de control de calidad} que mida la recurrencia del evento. Este patrón impide que los sistemas de salud aprendan de sus errores y escalen la vigilancia frente al volumen de notas clínicas generadas en hospitales de tercer nivel.

En consecuencia, el problema central se formula mediante la siguiente pregunta de investigación:

\begin{quote}
\textit{¿En qué medida un pipeline integrado de procesamiento de lenguaje natural y aprendizaje automático supervisado, aplicado sobre notas clínicas no estructuradas del \textit{dataset} MIMIC-IV, permite detectar eventos adversos y aplicar la Matriz de Priorización del Modelo GEMSES ---incluyendo la asignación automática del Nivel de Gestión y del responsable institucional--- con una concordancia significativamente mayor a la clasificación manual experta tradicional?}
\end{quote}

% ─────────────────────────────────────────────────────────────
\section{Justificación}
\label{sec:justificacion}

La presente investigación se justifica desde cuatro dimensiones complementarias:

% - - - - - - - - - - - - - - - - - - - - - - - - - - - - - -
\subsection{Justificación Teórica}
\label{subsec:justif_teorica}

El campo de la detección automática de eventos adversos en texto clínico ha avanzado sustancialmente en los últimos años, con trabajos basados en sistemas de supervisión débil (Ratner et al., 2017), modelos de transformers biomédicos (Lee et al., 2020; Alsentzer et al., 2019) y técnicas de negación clínica como NegEx (Chapman et al., 2001). Sin embargo, ningún trabajo previo ha articulado estos enfoques con la Matriz de Priorización del Modelo GEMSES ni con la taxonomía de 231 eventos del Anexo 02 de la Directiva GG-ESSALUD-2021. Esta investigación cierra esa brecha teórica, aportando un marco de operacionalización que vincula las dimensiones de Tiempo, Calidad, Costos y Satisfacción del Modelo GEMSES con variables derivables de datos clínicos estructurados y no estructurados. Los pesos de impacto por dimensión ya están precalculados en el Anexo~02 para cada uno de los 231 eventos; esta investigación los operacionaliza computacionalmente.

Adicionalmente, el estudio contribuye al debate metodológico sobre la transferibilidad de modelos de NLP entrenados en inglés clínico hacia sistemas sanitarios de habla hispana, problema de creciente relevancia en la literatura de informática biomédica internacional.

% - - - - - - - - - - - - - - - - - - - - - - - - - - - - - -
\subsection{Justificación Práctica}
\label{subsec:justif_practica}

Desde una perspectiva práctica, el sistema propuesto responde a una necesidad institucional concreta: el cuello de botella operativo que genera el procesamiento manual de notificaciones de eventos adversos en EsSalud. El \textit{pipeline} automatizado permitiría procesar el volumen completo de notas clínicas de alta (331{,}793 notas en MIMIC-IV; equivalente a decenas de miles en EsSalud) sin incremento proporcional de recursos humanos especializados, reduciendo el tiempo entre la ocurrencia del evento y su priorización para la gestión correctiva.

El modelo validado sobre MIMIC-IV generará un protocolo de implementación documentado y replicable, que podrá ser adaptado por la Oficina de Gestión de la Calidad y Humanización de EsSalud ---o cualquier institución que opere bajo estándares GEMSES--- cuando se habilite el acceso a datos clínicos institucionales para investigación. El impacto práctico esperado incluye: reducción de la subnotificación, mayor consistencia en la codificación, y escalamiento de la capacidad de vigilancia de la seguridad del paciente.

% - - - - - - - - - - - - - - - - - - - - - - - - - - - - - -
\subsection{Justificación Metodológica}
\label{subsec:justif_metodologica}

El uso de MIMIC-IV como corpus de validación garantiza la reproducibilidad del estudio: es un \textit{dataset} de acceso controlado pero abierto (PhysioNet \textit{credentialing}), con 145{,}914 pacientes únicos, en formato estándar internacional, y ampliamente utilizado como \textit{benchmark} en investigación de NLP clínico. La elección de este \textit{dataset} ---en lugar de datos institucionales peruanos no accesibles--- no es una limitación sino una decisión metodológica deliberada que garantiza la trazabilidad, la posibilidad de revisión por pares y la comparabilidad con el estado del arte internacional.

El diseño del \textit{pipeline} combina supervisión débil para etiquetado a escala, detección de negación con NegEx, mapeo a códigos ICD-10, y comparación de múltiples clasificadores (TF-IDF + regresión logística, LinearSVC, modelos basados en \textit{transformers}). La validación mediante panel experto ($n = 70$ notas) con métricas de concordancia (kappa de Cohen) proporciona evidencia estadística rigurosa sobre la capacidad del sistema, con un IC$_{95\%}$ Wilson calculado correctamente para muestras pequeñas.

% - - - - - - - - - - - - - - - - - - - - - - - - - - - - - -
\subsection{Justificación Social}
\label{subsec:justif_social}

La OMS estima que en países de ingresos medios y bajos, la carga de eventos adversos evitables equivale a 134 millones de episodios anuales con 2.6 millones de muertes asociadas (OMS, 2019). Perú, como país de ingreso medio-alto con un sistema de seguridad social en expansión, enfrenta este problema de manera aguda: la red asistencial de EsSalud atiende a más de 13 millones de asegurados con capacidad limitada de vigilancia activa de seguridad, que es detectada cuando el paciente denuncia; SUSALUD ha multado a instituciones por eventos de seguridad del paciente evitables por un valor de S/~8{,}482{,}906.50 en el año 2025.

Un sistema automatizado de detección y priorización contribuye directamente a la protección de pacientes vulnerables, a la rendición de cuentas institucional y al aprendizaje organizacional. La generación de alertas tempranas sobre eventos de nivel Rojo (alta prioridad GEMSES) permitiría activar protocolos de mitigación antes de que el daño escale, con un impacto social concreto y medible en términos de morbimortalidad prevenible y sostenibilidad financiera.


% Capítulo II: Objetivos
% ============================================================
%  CAPÍTULO II: OBJETIVOS
% ============================================================

\chapter{OBJETIVOS}
\label{cap:objetivos}

% ─────────────────────────────────────────────────────────────
\section{Objetivo General}
\label{sec:obj_general}

Diseñar, implementar y validar un sistema basado en procesamiento de lenguaje natural y aprendizaje automático que detecte eventos adversos a partir de notas clínicas no estructuradas del \textit{dataset} MIMIC-IV y los priorice automáticamente aplicando la Matriz de Priorización del Modelo GEMSES ---incluyendo el cálculo automatizado del Nivel de Gestión y la asignación del responsable institucional--- con miras a su transferencia metodológica a sistemas sanitarios de habla hispana y, en particular, a la red asistencial de EsSalud.

% ─────────────────────────────────────────────────────────────
\section{Objetivos Específicos}
\label{sec:obj_especificos}

Para alcanzar el objetivo general se plantean los siguientes objetivos específicos:

\begin{enumerate}[label=\textbf{OE\arabic*.}, leftmargin=*, itemsep=6pt]

    \item \textbf{Construir un corpus etiquetado} de eventos adversos a partir de notas clínicas no estructuradas de MIMIC-IV, utilizando como marco de etiquetado los 188 eventos detectables (de los 231 del Anexo~02 de la Directiva GG-ESSALUD-2021), clasificados por cuatro niveles de detección: (A)~ICD-10 + NLP, (B)~patrones de texto libre, (C)~variables estructuradas calculadas y (D)~auditoría de completitud de historia clínica.

    \item \textbf{Implementar y comparar} al menos tres modelos de procesamiento de lenguaje natural para la detección y clasificación de eventos adversos en texto libre, incluyendo modelos basados en \textit{transformers} biomédicos (BioBERT, ClinicalBERT) y un modelo base clásico de referencia (TF-IDF + regresión logística).

    \item \textbf{Automatizar la priorización GEMSES} aplicando, para cada evento detectado, la fórmula
    \[
        \text{Nivel de Riesgo} = \text{Frecuencia} \times \text{Impacto}
    \]
    usando los pesos de impacto ya establecidos en el Anexo~02 para cada uno de los 231 eventos (dimensiones: Estancia Hospitalaria, Complicaciones, Costos, Satisfacción del Paciente e Impacto Ambiental), y determinar automáticamente el Nivel de Gestión (Verde / Amarillo / Rojo) y el responsable institucional correspondiente.

    \item \textbf{Validar la precisión del \textit{pipeline}} comparando sus resultados contra la revisión manual de un panel de expertos sobre una muestra de 70 notas clínicas seleccionadas de forma estratificada, utilizando métricas estándar de concordancia e índices de clasificación para determinar si el sistema automatizado alcanza un nivel de acuerdo aceptable con el juicio experto.

    \item \textbf{Documentar el método de transferencia} del \textit{pipeline} a sistemas sanitarios de habla hispana, identificando los retos lingüísticos, taxonómicos y de localización que deberán abordarse para su implementación en EsSalud u otras instituciones equivalentes.

\end{enumerate}

\vspace{6pt}
Los objetivos OE1 y OE2 responden al obstáculo de la dependencia del análisis manual mediante la construcción de corpus etiquetado y la comparación de modelos. OE3 operacionaliza el instrumento GEMSES con sus pesos de impacto oficiales, automatizando la cadena completa Frecuencia $\times$ Impacto $\rightarrow$ Nivel de Gestión $\rightarrow$ Responsable. OE4 provee la evidencia empírica de validez con rigor estadístico. OE5 conecta el resultado de laboratorio con el impacto social y práctico, cerrando el ciclo de investigación aplicada y garantizando la coherencia global del proyecto.


% Los siguientes capítulos se agregarán conforme avance
% \include{3_1_CAPITULO_MARCO_TEORICO/Capitulo3}
% \include{4_1_CAPITULO_METODOLOGICO/Capitulo4}
% \include{5_1_CAPITULO_ADMIN/Capitulo5}

% ── PARTE FINAL ────────────────────────────────────────────
\backmatter

% ============================================================
%  REFERENCIAS BIBLIOGRÁFICAS
% ============================================================

\chapter*{REFERENCIAS BIBLIOGRÁFICAS}
\addcontentsline{toc}{chapter}{Referencias Bibliográficas}

% Completar con las referencias APA según avance la tesis.
% Ejemplo de entradas pendientes:

\begin{itemize}
    \item Alsentzer, E., Murphy, J. R., Boag, W., Weng, W.-H., Jin, D., Naumann, T., \& McDermott, M. (2019). Publicly available clinical BERT embeddings. \textit{arXiv preprint arXiv:1904.03323}.

    \item Chapman, W. W., Bridewell, W., Hanbury, P., Cooper, G. F., \& Buchanan, B. G. (2001). A simple algorithm for identifying negated findings and diseases in discharge summaries. \textit{Journal of Biomedical Informatics, 34}(5), 301--310.

    \item Lee, J., Yoon, W., Kim, S., Kim, D., Kim, S., So, C. H., \& Kang, J. (2020). BioBERT: a pre-trained biomedical language representation model for biomedical text mining. \textit{Bioinformatics, 36}(4), 1234--1240.

    \item Organización Mundial de la Salud. (2019). \textit{Acción mundial para la seguridad del paciente}. OMS.

    \item Ratner, A., De Sa, C., Wu, S., Selsam, D., \& Ré, C. (2017). Data programming: Creating large training sets, quickly. \textit{Advances in Neural Information Processing Systems, 29}.
\end{itemize}

% ============================================================
%  ANEXOS
% ============================================================

\chapter*{ANEXOS}
\addcontentsline{toc}{chapter}{Anexos}

% Pendiente: incluir Matriz de Consistencia y demás anexos
% conforme avanza la investigación.


\end{document}
